\subsection{The Explanatory Chain and the Content of the Programme}
  \label{subsec:time-irreversibility-and-cosmological-expansion}

  The objective and the foundational requirement of the preceding subsections fix
  what the framework must deliver.
  If established physical theories are effective descriptions, accounting for them
  requires more than placing them inside a broader interpretation.
  It requires an explicit chain of consequences showing how their characteristic
  structures become necessary from the admissibility primitive, while preserving
  their validated empirical domains.
  This explanatory chain constitutes the principal content of the programme.
  Its purpose is not to replace successful theories by verbal reinterpretations, but
  to identify which of their independent inputs can be recovered as structural
  outputs, which results remain conditional on additional bridges, and which claims
  remain open.
  Each link must therefore state both the structure obtained and its epistemic
  status.
  Starting from the local structure of admissible non-injective transitions, the
  explananda form an ordered programme: temporal ordering and irreversibility;
  effective spatial separation and geometry; Lorentzian causal structure; quantum
  states, coherent amplitudes, the Born rule, and Bell non-factorisability; internal
  gauge structure; fermionic matter and generation structure; gauge and
  gravitational dynamics; and cosmological evolution.
  The order is logical rather than merely expository: later sectors may use
  structures established by earlier ones, and no downstream result should be
  imported into the derivation of its own prerequisites.
  The first part of the programme is therefore reconstructive.
  It seeks to derive structures that established theories introduce independently
  and to determine the precise domain in which the resulting effective descriptions
  apply.
  The second part is prospective: the observational signatures collected in
  Section~\ref{sec:testable-predictions-and-observational-signatures} provide the
  means by which the framework can be distinguished, constrained, or rejected.
  The first link in this explanatory chain is temporal ordering, together with the
  irreversibility and cosmological behaviour that follow from it.

  In standard cosmology, expansion is described kinematically through the scale
  factor, while the arrow of time is typically attributed to boundary conditions
  or entropy
  growth~\cite{peebles1993principles,Prigogine1997,penrose1989emperors}.
  Neither account addresses the structural origin of temporal ordering itself.

  In Cosmochrony, temporal ordering is not postulated as a property of the
  substrate~$\chi$, nor derived from an underlying dynamics.
  It emerges instead as an induced feature of the succession of admissible
  transitions: the cumulative projected spectral capacity~$I(n)$, measuring the
  spectral structure stably realised at each step of the admissibility cascade, is
  monotone non-decreasing along all tested admissible trajectories~\cite{Beau2026pto}.
  This monotonicity defines an intrinsic ordering on effective states without
  requiring any evolution of~$\chi$ itself.
  One-way activation is the condition of possibility for this ordering: because
  spectral modes, once stably projected, do not deactivate, $I(n)$ is
  well-defined as a globally consistent monotone quantity, and irreversibility
  follows.
  One-way activation is a property of \emph{spectral modes}, not of the
  coherence of projected configurations; the structural reorganization of
  unstable configurations into more stable ones is separately accounted for
  by the admissible factorization mechanism
  (Section~\ref{subsec:metastability-and-decay}).

  Three levels of status attach to this picture and should be kept distinct.
  The monotonicity of~$I(n)$ is an \emph{established numerical result}, confirmed
  across~$169$ conjugate pairs and~$310$ post-burn-in BFS steps for
  $q \in \{29, 61, 101, 151\}$, with zero physically meaningful
  regression~\cite{Beau2026pto}.
  The one-way activation property that underpins it --- formally, that
  $\sigma_{\mathrm{pair}}(n) \le \sigma_{\mathrm{pair}}(n-1)$ for all
  admissible~$q$, all conjugate pairs, and all $n \ge 1$ --- is an
  \emph{open conjecture}: its analytical proof from the Gram--Schmidt structure,
  the Born--Infeld admissibility constraints, and the non-injectivity of~$\Pi$
  remains to be established
  (see Section~\ref{subsec:monotonicity-and-arrow-of-time}).
  The identification of this ordering with physical time is an
  \emph{interpretation}: temporal ordering is the relational structure induced
  by admissible projections, not an evolution of~$\chi$ itself.

  Cosmological expansion admits an analogous interpretation, distinct in status
  from the ordering result above.
  The progressive growth of~$I(n)$ at the global scale corresponds to an ever
  broader range of mutually distinguishable projected configurations becoming
  stably realised.
  In this interpretation, expansion is not driven by an external energy component
  but reflects the monotone accumulation of admissible projected structure.
  This identification is interpretive: it does not follow from non-injectivity
  alone, and the exact links and their statuses are addressed separately in the
  cosmological sector.
  In the effective geometric description, this global growth manifests as the
  large-scale expansion of the observable universe.
